\section{Hawking 奇性定理}\label{sec:hawking-singularity-theorem}

如第~\ref{sec:schwarzschild-solution}~节和第~\ref{sec:cosmology}~节所见, Schwarzschild 解和 FLRW 宇宙学模型都呈现奇点, 类时测地线和零模测地线越过这些奇点后无法继续延拓.

\begin{RiemannDefinition}\label{def:singular-spacetime}
若时空 $(M,g)$ 不是测地完备的, 则称其为\textbf{奇性时空}.
\end{RiemannDefinition}

一度有人认为, 上述例子的奇性是由它们高度对称造成的, 更真实的时空将没有奇性. 沿着 Hawking 和 Penrose 的思路 [\RGref{bib:Pen65}, \RGref{bib:Haw67}, \RGref{bib:HP70}], 我们将证明事实并非如此: 对这些解作任何足够小的扰动, 所得时空仍有奇性.

给定 Riemann 流形是否测地完备的问题, 由 Hopf–Rinow 定理判定. 遗憾的是, 这一结论在 Lorentz 几何中不成立 (根本原因是无法用度量定义距离函数). 例如, 紧流形不一定测地完备 [参见 Exercise~\ref{ex:hawking-singularity-theorem}(\ref{ex:clifton-pohl-torus-incompleteness})], 而测地完备流形上的指数映射也不一定是满射 [参见 Exercise~\ref{ex:hawking-singularity-theorem}(\ref{ex:anti-de-sitter-geodesic-completeness})].

设 $(M,g)$ 是整体双曲时空, $S$ 是带有指向未来的单位法向量场 $n$ 的 Cauchy 超曲面. 对每个 $p\in S$, 设 $c_p$ 为初始条件为 $n_p$ 的类时测地线. 在包含 $\{0\}\times S$ 的开集 $U\subset\mathbb{R}\times S$ 上, 定义光滑映射 $\exp:U\to M$, 使 $\exp(t,p)=c_p(t)$.

\begin{RiemannDefinition}\label{def:conjugate-points-to-cauchy-hypersurface}
$\exp$ 的临界值称为 $S$ 的\textbf{共轭点}.
\end{RiemannDefinition}

粗略地说, 共轭点是从 $S$ 上邻近点正交出发的测地线相交之处.

设 $q=\exp(t_0,p)$ 是不与 $S$ 共轭的一点, $(x^1,x^2,x^3)$ 是 $S$ 上 $p$ 附近的局部坐标. 则 $(t,x^1,x^2,x^3)$ 是某个开集 $V\ni q$ 上的局部坐标. 由于 $\frac{\partial}{\partial t}$ 是与 $S$ 正交的测地线的单位切向量场, 有 $g_{00}=\langle\frac{\partial}{\partial t},\frac{\partial}{\partial t}\rangle=-1$. 另一方面,
\begin{align*}
  \frac{\partial g_{0i}}{\partial t}
   &=\frac{\partial}{\partial t}\left\langle\frac{\partial}{\partial t},
              \frac{\partial}{\partial x^i}\right\rangle
    =\left\langle\frac{\partial}{\partial t},
              \nabla_{\frac{\partial}{\partial t}}\frac{\partial}{\partial x^i}\right\rangle\\
   &=\left\langle\frac{\partial}{\partial t},
              \nabla_{\frac{\partial}{\partial x^i}}\frac{\partial}{\partial t}\right\rangle
    =\frac{1}{2}\frac{\partial}{\partial x^i}
              \left\langle\frac{\partial}{\partial t},\frac{\partial}{\partial t}\right\rangle=0
\end{align*}
对 $i=1,2,3$ 成立, 而且由于在 $S$ 上 $g_{0i}=0$, 在 $V$ 上也有 $g_{0i}=0$. 因此 $t$ 为常数的曲面与切向量为 $\frac{\partial}{\partial t}$ 的测地线正交. 为此, 称 $(t,x^1,x^2,x^3)$ 为\textbf{同步坐标系}. 在此坐标系下有
\[
  g=-\mathrm{d}t\otimes\mathrm{d}t
      +\sum_{i,j=1}^3\gamma_{ij}\,\mathrm{d}x^i\otimes\mathrm{d}x^j,
\]
其中函数
\[
  \gamma_{ij}:=\left\langle\frac{\partial}{\partial x^i},
                                  \frac{\partial}{\partial x^j}\right\rangle
\]
组成一个正定矩阵. 由于向量场 $\frac{\partial}{\partial x^i}$ 总能沿 $c_p$ 定义, 即便在同步坐标系失效的点, 即在 $S$ 的共轭点处, 矩阵 $(\gamma_{ij})$ 也沿 $c_p$ 良定义. 这些正是 $\gamma(t):=\det(\gamma_{ij}(t))$ 为零的点, 因为只有在这些点处, $\{\frac{\partial}{\partial t},\frac{\partial}{\partial x^1},\frac{\partial}{\partial x^2},\frac{\partial}{\partial x^3}\}$ 才不再线性无关. (事实上, 向量场 $\frac{\partial}{\partial x^i}$ 是沿 $c_p$ 的 Jacobi 向量场 --- 见第~\ref{chap:riemannian-manifolds}~章的 Exercise~\ref{ex:minimizing-properties-of-geodesics}(\ref{ex:jacobi-fields-and-cut-locus}).)

容易看出
\[
  \Gamma_{00}^0=\Gamma_{00}^i=0
  \qquad\text{且}\qquad
  \Gamma_{0j}^i=\sum_{k=1}^3\gamma^{ik}\beta_{kj},
\]
其中 $(\gamma^{ij})=(\gamma_{ij})^{-1}$, $\beta_{ij}=\frac{1}{2}\frac{\partial\gamma_{ij}}{\partial t}$ [参见 Exercise~\ref{ex:hawking-singularity-theorem}(\ref{ex:christoffel-symbols-in-synchronized-coordinates})]. 因而
\begin{align*}
  R_{00}
  &=\sum_{i=1}^3R_{i00}^{\ \ \ i}
   =\sum_{i=1}^3\left(\frac{\partial\Gamma_{00}^i}{\partial x^i}
     -\frac{\partial\Gamma_{i0}^i}{\partial t}
     +\sum_{j=1}^3\Gamma_{00}^j\Gamma_{ij}^i
     -\sum_{j=1}^3\Gamma_{i0}^j\Gamma_{0j}^i\right)\\
  &=-\frac{\partial}{\partial t}\left(\sum_{i,j=1}^3\gamma^{ij}\beta_{ij}\right)
     -\sum_{i,j,k,l=1}^3\gamma^{jk}\gamma^{il}\beta_{ki}\beta_{lj}.
\end{align*}
(参见第~\ref{chap:curvature}~章第~\ref{subsec:notes-manifolds-of-constant-curvature}~节.) 此式中出现的量
\[
  \theta:=\sum_{i,j=1}^3\gamma^{ij}\beta_{ij}
\]
称为同步观察者的\textbf{膨胀率}, 具有重要的几何意义:
\[
  \theta=\frac{1}{2}\operatorname{tr}\left((\gamma_{ij})^{-1}
      \frac{\partial}{\partial t}(\gamma_{ij})\right)
  =\frac{1}{2}\frac{\partial}{\partial t}\log\gamma
  =\frac{\partial}{\partial t}\log\gamma^{1/2}.
\]
这里用了对任意光滑矩阵函数 $A:\mathbb{R}\to GL(n)$ 成立的公式
\[
  (\log(\det A))'=\operatorname{tr}(A^{-1}A')
\]
[参见第~\ref{chap:differentiable-manifolds}~章的 Example~\ref{ex:classical-matrix-lie-groups}(\ref{ex:special-linear-group-regular-level-set})]. 因此, 膨胀率给出了同步观察者测得的三维体积元的变化. 对我们而言更重要的是, 膨胀率的奇点表明 $\gamma$ 取零值, 即出现 $S$ 的一个共轭点.

\begin{RiemannDefinition}\label{def:strong-energy-condition}
若对任意类时向量场 $V\in\mathfrak{X}(M)$ 均有 $\operatorname{Ric}(V,V)\geq0$, 则称时空 $(M,g)$ 满足\textbf{强能量条件}.
\end{RiemannDefinition}

由 Einstein 方程, 这等价于要求对任意类时向量场 $V\in\mathfrak{X}(M)$, 约化能量--动量张量 $T$ 满足 $T(V,V)\geq0$. 对于静止质量密度函数为 $\rho\in C^\infty(M)$, 单位速度向量场为 $U\in\mathfrak{X}(M)$ 的无压流体, 此要求成为
\[
  \rho\left(\langle U,V\rangle^2+\frac{1}{2}\langle V,V\rangle\right)\geq0,
\]
或者, 由于括号内的项总为正 [参见 Exercise~\ref{ex:hawking-singularity-theorem}(\ref{ex:positive-pressureless-fluid-energy-contraction})], 即 $\rho\geq0$. 对于更复杂的物质模型, 强能量条件也给出同样合理的限制.

\begin{RiemannProposition}\label{prop:timelike-focusing-under-strong-energy-condition}
设 $(M,g)$ 为满足强能量条件的整体双曲时空, $S\subset M$ 为 Cauchy 超曲面, $p\in S$ 为使 $\theta=\theta_0<0$ 的一点. 则测地线 $c_p$ 至少包含一个 $S$ 的共轭点, 距 $S$ 向未来至多为 $-\frac{3}{\theta_0}$ (假定它能延拓至该距离).
\end{RiemannProposition}

\begin{RiemannProof}
由于 $(M,g)$ 满足强能量条件, 在任何同步标架上都有 $R_{00}=\operatorname{Ric}(\frac{\partial}{\partial t},\frac{\partial}{\partial t})\geq0$. 因此在这样的标架上
\[
  \frac{\partial\theta}{\partial t}
  +\sum_{i,j,k,l=1}^3\gamma^{jk}\gamma^{il}\beta_{ki}\beta_{lj}\leq0.
\]
选取一个标准正交基 (此时 $\gamma^{ij}=\delta_{ij}$), 利用对 $n\times n$ 方阵成立的不等式
\[
  (\operatorname{tr}A)^2\leq n\operatorname{tr}(A^tA),
\]
(它是 Cauchy–Schwarz 不等式的直接推论), 容易证明
\[
  \sum_{i,j,k,l=1}^3\gamma^{jk}\gamma^{il}\beta_{ki}\beta_{lj}
   =\sum_{i,j=1}^3\beta_{ji}\beta_{ij}
   =\operatorname{tr}\bigl((\beta_{ij})(\beta_{ij})^t\bigr)\geq\frac{1}{3}\theta^2.
\]
所以 $\theta$ 必定满足
\[
  \frac{\partial\theta}{\partial t}+\frac{1}{3}\theta^2\leq0.
\]
对这一不等式积分得到
\[
  \frac{1}{\theta}\geq\frac{1}{\theta_0}+\frac{t}{3},
\]
因此 $\theta$ 必定在 $t$ 不大于 $-\frac{3}{\theta_0}$ 时发散.
\end{RiemannProof}

\begin{RiemannProposition}\label{prop:timelike-geodesic-not-maximal-past-conjugate-point}
设 $(M,g)$ 为整体双曲时空, $S$ 为 Cauchy 超曲面, $p\in M$, $c$ 为过 $p$ 且与 $S$ 正交的类时测地线. 若在 $S$ 与 $p$ 之间存在一个共轭点, 则 $c$ (在连接 $S$ 与 $p$ 的类时曲线中) 不使长度达到最大.
\end{RiemannProposition}

\begin{RiemannProof}
只给出证明的概要. 设 $q$ 为沿 $c$ 在 $S$ 与 $p$ 之间的第一个共轭点. 可在 $c$ 从 $S$ 到 $q$ 的部分周围使用同步坐标系. 由于 $q$ 与 $S$ 共轭, 存在另一条正交于 $S$, 长度同为 $t(q)$ 且在 $q$ 处 (近似) 与 $c$ 相交的测地线 $\widetilde c$. 设 $V$ 是 $q$ 的测地凸邻域, $r\in V$ 为 $\widetilde c$ 上介于 $S$ 与 $q$ 之间的一点, $s\in V$ 为 $c$ 上介于 $q$ 与 $p$ 之间的一点 (参见 Fig.~\ref{fig:nonmaximal-geodesic-past-conjugate-point}). 先沿 $\widetilde c$ 从 $S$ 到 $r$, 再沿 $V$ 中唯一连接 $r$ 与 $s$ 的测地线, 最后沿 $c$ 从 $s$ 到 $p$, 所得分段光滑的类时曲线连接 $S$ 与 $p$, 且长度严格大于 $c$ 的长度 (由广义双生子佯谬). 容易将这条曲线光滑化, 同时保持其长度大于 $c$ 的长度.
\end{RiemannProof}
\RGfigure[.60\textwidth]{nonmaximal-geodesic-past-conjugate-point}{Proposition~\ref{prop:timelike-geodesic-not-maximal-past-conjugate-point} 的证明}

\begin{RiemannProposition}\label{prop:compact-past-intersect-future-domain-of-dependence}
设 $(M,g)$ 为整体双曲时空, $S$ 为 Cauchy 超曲面, 且 $p\in D^+(S)$. 则 $D^+(S)\cap J^-(p)$ 是紧的.
\end{RiemannProposition}

\begin{RiemannProof}
定义\textbf{简单邻域} $U\subset M$ 为与开球微分同胚的测地凸开集, 它的边界是一个更大的测地凸开集的紧子流形 (因此 $\partial U$ 与 $S^3$ 微分同胚, 且 $\overline U$ 是紧的). 显然, 简单邻域构成 $M$ 的拓扑的一组基. 此外容易证明, 任意开覆盖 $\{V_\alpha\}_{\alpha\in A}$ 都有由简单邻域组成的可数局部有限加细 $\{U_n\}_{n\in\mathbb{N}}$ [参见 Exercise~\ref{ex:hawking-singularity-theorem}(\ref{ex:locally-finite-refinements-by-simple-neighborhoods})].

若 $A=D^+(S)\cap J^-(p)$ 不是紧的, 则存在 $A$ 的一个由简单邻域组成的可数局部有限开覆盖 $\{U_n\}_{n\in\mathbb{N}}$, 它没有有限子覆盖. 取 $q_n\in A\cap U_n$, 使当 $m\ne n$ 时 $q_m\ne q_n$. 序列 $\{q_n\}_{n\in\mathbb{N}}$ 不可能有聚点, 因为 $M$ 中任一点都有一个只与有限个简单邻域 $U_n$ 相交的邻域. 特别地, 每个简单邻域 $U_n$ 只含序列中有限多个点 (因为 $\overline U_n$ 紧).

令 $p_1=p$. 由于 $p_1\in A$, 存在某个 $n_1\in\mathbb{N}$ 使 $p_1\in U_{n_1}$. 设 $q_n\notin U_{n_1}$. 由于 $q_n\in J^-(p_1)$, 存在朝向未来的因果曲线 $c_n$ 连接 $q_n$ 至 $p_1$. 这条曲线必与 $\partial U_{n_1}$ 相交. 设 $r_{1,n}$ 为一个交点. 由于 $U_{n_1}$ 只含序列 $\{q_n\}_{n\in\mathbb{N}}$ 中有限多个点, 将有无穷多个交点 $r_{1,n}$. 由于 $\partial U_{n_1}$ 紧, 它们将聚集到某点 $p_2\in\partial U_{n_1}$ (参见 Fig.~\ref{fig:compact-causal-past-limit-curve}).

由于 $\overline U_{n_1}$ 包含于测地凸开集 $V$, 而 $V$ 可选得使 $v\mapsto(\pi(v),\exp(v))$ 为到 $V\times V$ 的微分同胚, 所以有 $p_2\in J^-(p_1)$: 若 $\gamma_{1,n}$ 为连接 $p_1$ 至 $r_{1,n}$ 的唯一因果测地线, 以全局时间函数 $t:M\to\mathbb{R}$ 参数化, 则 $\{\gamma_{1,n}\}$ 中与 $\{r_{1,n}\}$ 的某个收敛子序列对应的子序列, 收敛到连接 $p_1$ 至 $p_2$ 的因果测地线 $\gamma_1$. 若 $S=t^{-1}(0)$, 则有 $t(r_{1,n})\geq0$, 蕴含 $t(p_2)\geq0$, 因而 $p_2\in A$. 由于 $p_2\notin U_{n_1}$, 必存在 $n_2\in\mathbb{N}$ 使 $p_2\in U_{n_2}$.

由于 $U_{n_2}$ 只含序列 $\{q_n\}_{n\in\mathbb{N}}$ 中有限多个点, 无穷多条曲线 $c_n$ 必在 $r_{1,n}$ 的过去与 $\partial U_{n_2}$ 相交. 设 $r_{2,n}$ 为这些交点. 由于 $\partial U_{n_2}$ 紧, $\{r_{2,n}\}$ 必聚集到某点 $p_3\in\partial U_{n_2}$. 由于 $\overline U_{n_2}$ 包含于一个测地凸开集, 有 $p_3\in J^-(p_2)$: 若 $\gamma_{2,n}$ 为连接 $r_{1,n}$ 至 $r_{2,n}$ 的唯一因果测地线, 以全局时间函数参数化, 则 $\{\gamma_{2,n}\}$ 中与 $\{r_{1,n}\}$ 和 $\{r_{2,n}\}$ 的收敛子序列对应的子序列, 收敛到连接 $p_2$ 至 $p_3$ 的因果测地线. 由于 $J^-(p_2)\subset J^-(p_1)$, 且 $t(r_{2,n})\geq0\implies t(p_3)\geq0$, 有 $p_3\in A$.

重复以上过程, 可构造 $A$ 中的点列 $\{p_i\}_{i\in\mathbb{N}}$, 满足 $p_i\in U_{n_i}$, 当 $i\ne j$ 时 $n_i\ne n_j$, 且 $p_i$ 由因果测地线 $\gamma_i$ 连接到 $p_{i+1}$. 显然, $\gamma_i$ 不可能与 $S$ 相交, 因为 $t(p_{i+1})>t(p_{i+2})\geq0$. 另一方面, 连接各曲线 $\gamma_i$ 所得的分段光滑因果曲线, 容易光滑化为从 $p_1$ 出发且不与 $S$ 相交的朝向过去的因果曲线. 最后, 此曲线不可延拓: 由于 $\{p_i\}_{i\in\mathbb{N}}$ 不可能聚集, 它不可能收敛到任何一点. 但由于 $p_1\in D^+(S)$, 这条曲线必须与 $S$ 相交. 因此 $A$ 必为紧集.
\end{RiemannProof}
\RGfigure[.60\textwidth]{compact-causal-past-limit-curve}{Proposition~\ref{prop:compact-past-intersect-future-domain-of-dependence} 的证明}

\begin{RiemannCorollary}\label{cor:closed-causal-futures-and-compact-diamonds}
设 $(M,g)$ 为整体双曲时空, $p,q\in M$. 则:
\begin{enumerate}[label=(\roman*),leftmargin=*]
\item $J^+(p)$ 是闭集;
\item $J^+(p)\cap J^-(q)$ 是紧集.
\end{enumerate}
\end{RiemannCorollary}

\begin{RiemannProof}
Exercise~\ref{ex:hawking-singularity-theorem}(\ref{ex:prove-closed-causal-futures-and-compact-diamonds}).
\end{RiemannProof}

Proposition~\ref{prop:compact-past-intersect-future-domain-of-dependence} 是建立下列基本结果的关键.

\begin{RiemannTheorem}\label{thm:existence-of-maximal-timelike-geodesic-from-cauchy-hypersurface}
设 $(M,g)$ 为具有 Cauchy 超曲面 $S$ 的整体双曲时空, 且 $p\in D^+(S)$. 则在所有连接 $p$ 至 $S$ 的类时曲线中, 存在一条长度最大的类时曲线. 这条曲线是正交于 $S$ 的类时测地线.
\end{RiemannTheorem}

\begin{RiemannProof}
考虑所有连接 $S$ 至 $p$ 的类时曲线所成的集合 $T(S,p)$. 由于总能以全局时间函数 $t:M\to\mathbb{R}$ 为参数, 这些曲线由它们的像确定, 而这些像是紧集 $A=D^+(S)\cap J^-(p)$ 的紧子集. 众所周知 (参见 \RGcite{ch06-Mun00}), $A$ 的所有紧子集所成的集合 $C(A)$ 在 \textbf{Hausdorff 度量} $d_H$ 下是紧度量空间. 此度量定义如下: 若 $d:M\times M\to\mathbb{R}$ 是给出 $M$ 的拓扑的度量, 则
\[
  d_H(K,L)=\inf\{\varepsilon>0\mid K\subset U_\varepsilon(L)
       \ \text{且}\ L\subset U_\varepsilon(K)\},
\]
其中 $U_\varepsilon(K)$ 为 $K$ 关于度量 $d$ 的 $\varepsilon$-邻域. 因此, 闭包 $C(S,p):=\overline{T(S,p)}$ 是 $C(A)$ 的紧子集. 不难证明, $C(S,p)$ 可以与连接 $S$ 至 $p$ 的\textbf{连续因果曲线}的集合等同 (若连续曲线 $c:[0,t(p)]\to M$ 在 $t_2>t_1$ 时都有 $c(t_2)\in J^+(c(t_1))$, 则称其为\textbf{因果曲线}).

长度函数 $\tau:T(S,p)\to\mathbb{R}$ 定义为
\[
  \tau(c):=\int_0^{t(p)}|\dot c(t)|\,\mathrm{d}t.
\]
这个函数是\phantomsection\label{idx:hausdorff-upper-semicontinuity}\textbf{上半连续的}, 即对于 $\mathbb{R}$ 上的拓扑
\[
  \mathcal{O}=\{(-\infty,a)\mid-\infty\leq a\leq+\infty\}
\]
连续. 事实上, 设 $c\in T(S,p)$ 以其弧长 $u$ 参数化. 对于足够小的 $\varepsilon>0$, 函数 $u$ 可以延拓至 $\varepsilon$-邻域 $U_\varepsilon(c)$, 使其水平超曲面为类空的且与 $c$ 正交, 即 $-\operatorname{grad}u$ 是类时向量场, 并且在 $c$ 上与 $\dot c$ 一致 (参见 Fig.~\ref{fig:hausdorff-neighborhood-of-timelike-curve}). 若 $\gamma\in T(S,p)$ 位于 Hausdorff 度量 $d_H$ 的开球 $B_\varepsilon(c)\subset C(A)$ 中, 则可用 $u$ 作参数, 得到
\[
  \mathrm{d}u(\dot\gamma)=1\iff\langle\dot\gamma,\operatorname{grad}u\rangle=1.
\]
所以 $\dot\gamma$ 可分解为
\[
  \dot\gamma=\frac{1}{\langle\operatorname{grad}u,\operatorname{grad}u\rangle}
                       \operatorname{grad}u+X,
\]
其中 $X$ 为类空向量, 并与 $\operatorname{grad}u$ 正交, 因而
\[
  |\dot\gamma|=
  \left|\frac{1}{\langle\operatorname{grad}u,\operatorname{grad}u\rangle}
                                         +\langle X,X\rangle\right|^{1/2}.
\]
给定 $\delta>0$, 可以选取足够小的 $\varepsilon>0$, 使在 $\varepsilon$-邻域 $U_\varepsilon(c)$ 上有
\[
  -\frac{1}{\langle\operatorname{grad}u,\operatorname{grad}u\rangle}
     <\left(1+\frac{\delta}{2\tau(c)}\right)^2
\]
(因为沿 $c$ 有 $\langle\operatorname{grad}u,\operatorname{grad}u\rangle=-1$). 有
\[
  \tau(\gamma)=\int_0^{t(p)}\left|\frac{\mathrm{d}\gamma}{\mathrm{d}t}\right|\,\mathrm{d}t
  =\int_0^{t(p)}|\dot\gamma|\frac{\mathrm{d}u}{\mathrm{d}t}\,\mathrm{d}t
  =\int_{u(\gamma\cap S)}^{\tau(c)}|\dot\gamma|\,\mathrm{d}u,
\]
其中必须考虑 $c$ 不一定与 $S$ 正交, 因而 $\gamma$ 的起点不一定在 $u=0$ 处 (参见 Fig.~\ref{fig:hausdorff-neighborhood-of-timelike-curve}). 因此,
\begin{align*}
  \tau(\gamma)
  &=\int_{u(\gamma\cap S)}^{\tau(c)}
     \left|-\frac{1}{\langle\operatorname{grad}u,\operatorname{grad}u\rangle}
                                        -\langle X,X\rangle\right|^{1/2}\,\mathrm{d}u\\
  &<\int_{u(\gamma\cap S)}^{\tau(c)}\left(1+\frac{\delta}{2\tau(c)}\right)\,\mathrm{d}u
    =\left(1+\frac{\delta}{2\tau(c)}\right)(\tau(c)-u(\gamma\cap S)).
\end{align*}
取 $\varepsilon$ 足够小, 使在 $S\cap U_\varepsilon(c)$ 上有
\[
  |u|<\left(\frac{1}{\tau(c)}+\frac{2}{\delta}\right)^{-1},
\]
得到 $\tau(\gamma)<\tau(c)+\delta$, 这就证明了 $\tau$ 在 $T(S,p)$ 上的上半连续性. 因此, 长度函数可通过
\[
  \tau(c)=\lim_{\varepsilon\to0}\sup\{\tau(\gamma)\mid\gamma\in B_\varepsilon(c)\cap T(S,p)\}
\]
延拓至 $C(S,p)$ (因为 $\varepsilon>0$ 足够小时上确界有限). 此外, 显然若 $c\in T(S,p)$, 长度的上半连续性使 $\tau(c)$ 的两个定义一致. 长度函数到 $C(S,p)$ 的延拓显然是上半连续的: 给定 $c\in C(S,p)$ 及 $\delta>0$, 取 $\varepsilon>0$, 使对任意 $\gamma\in B_{2\varepsilon}(c)\cap T(S,p)$, 有 $\tau(\gamma)<\tau(c)+\frac{\delta}{2}$. 则显然对任意 $c'\in B_\varepsilon(c)$, 都有 $\tau(c')\leq\tau(c)+\frac{\delta}{2}<\tau(c)+\delta$.

最后, 注意 $\mathbb{R}$ 在拓扑 $\mathcal{O}$ 下的紧集就是具有最大元的集合. 因此, 长度函数在某点 $c\in C(S,p)$ 取得最大值. 只需再证明最大值也在光滑类时曲线 $\gamma$ 处取得. 为此, 用有限个测地凸邻域覆盖 $c$, 在 $c$ 上取点 $p_1,\ldots,p_k$, 使 $p_1\in S$, $p_k=p$, 并且对所有 $i=2,\ldots,k$, $c$ 在 $p_{i-1}$ 和 $p_i$ 之间的部分包含于一个测地凸邻域. 显然存在序列 $c_n\in T(S,p)$, 使 $c_n\to c$ 且 $\tau(c_n)\to\tau(c)$. 令 $t_i=t(p_i)$, $p_{i,n}$ 为 $c_n$ 与 $t^{-1}(t_i)$ 的交点. 以分段测地曲线 $\gamma_n$ 代替 $c_n$, 其中 $\gamma_n$ 由在相应的测地凸邻域中连接 $p_{i-1,n}$ 至 $p_{i,n}$ 而得. 则 $\tau(\gamma_n)\geq\tau(c_n)$, 从而 $\tau(\gamma_n)\to\tau(c)$. 因每个序列 $p_{i,n}$ 都收敛至 $p_i$, $\gamma_n$ 收敛至由连接 $p_{i-1}$ 至 $p_i$ ($i=2,\ldots,k$) 所得的分段测地曲线 $\gamma$, 且显然 $\tau(\gamma_n)\to\tau(\gamma)=\tau(c)$. 因而 $\gamma$ 是长度取最大值的一点. 最后, 注意 $\gamma$ 在各点 $p_i$ 必光滑, 否则可以利用广义双生子佯谬增加其长度. 因此 $\gamma$ 必为类时测地线. 利用 $\gamma(0)$ 附近的同步坐标系, 可见 $\gamma$ 必与 $S$ 正交, 否则就能增加其长度.
\end{RiemannProof}
\RGfigure[.60\textwidth]{hausdorff-neighborhood-of-timelike-curve}{Theorem~\ref{thm:existence-of-maximal-timelike-geodesic-from-cauchy-hypersurface} 的证明}

现在已经具备证明 Hawking 奇性定理的所有条件:

\begin{RiemannTheorem}[Hawking]\label{thm:hawking-timelike-singularity}
设 $(M,g)$ 为满足强能量条件的整体双曲时空, 并假设在一个 Cauchy 超曲面 $S$ 上, 膨胀率满足 $\theta\leq\theta_0<0$. 则 $(M,g)$ 是奇性时空.
\end{RiemannTheorem}

\begin{RiemannProof}
将证明与 $S$ 正交的任何朝向未来的类时测地线都不能在 $S$ 的未来延拓到大于 $\tau_0=-\frac{3}{\theta_0}$ 的固有时. 假设不然. 则对某个 $\varepsilon>0$, 存在一条定义在区间 $[0,\tau_0+\varepsilon]$ 上, 与 $S$ 正交且以固有时参数化的朝向未来的类时测地线 $c$. 令 $p=c(\tau_0+\varepsilon)$. 由 Theorem~\ref{thm:existence-of-maximal-timelike-geodesic-from-cauchy-hypersurface}, 存在连接 $S$ 至 $p$, 与 $S$ 正交且长度最大的类时测地线 $\gamma$. 因为 $\tau(c)=\tau_0+\varepsilon$, 必有 $\tau(\gamma)\geq\tau_0+\varepsilon$. Proposition~\ref{prop:timelike-focusing-under-strong-energy-condition} 保证 $\gamma$ 在 $S$ 的未来至多距离 $\tau_0$ 处出现一共轭点, 而 Proposition~\ref{prop:timelike-geodesic-not-maximal-past-conjugate-point} 断言 $\gamma$ 越过该点后不再使长度最大. 这就得出矛盾.
\end{RiemannProof}

\begin{RiemannRemark}\label{rem:past-timelike-incompleteness-from-positive-expansion}
显然, 若把 Cauchy 超曲面 $S$ 上的条件 $\theta\leq\theta_0<0$ 改为 $S$ 上的条件 $\theta\geq\theta_0>0$, 则 $(M,g)$ 仍是奇性时空. 在此情形, 与 $S$ 正交的任何\textbf{朝向过去}的类时测地线, 都不能在 $S$ 的\textbf{过去}延拓到大于 $\tau_0=\frac{3}{\theta_0}$ 的固有时.
\end{RiemannRemark}

\begin{RiemannExample}\label{ex:hawking-singularities-in-flrw-and-schwarzschild}
\begin{enumerate}[label=(\arabic*),ref=\arabic*,leftmargin=*]
\item FLRW 模型整体双曲 [参见 Exercise~\ref{ex:causality}(\ref{ex:globally-hyperbolic-spacetime-examples})], 且满足强能量条件 (因为 $\rho>0$). 此外,
\[
  \beta_{ij}=\frac{\dot a}{a}\gamma_{ij}\implies\theta=\frac{3\dot a}{a}.
\]
假定模型在时刻 $t_0$ 膨胀. 则在 Cauchy 超曲面 $S=\{t=t_0\}$ 上有 $\theta=\theta_0=\frac{3\dot a(t_0)}{a(t_0)}>0$, 因而 Theorem~\ref{thm:hawking-timelike-singularity} 保证模型在 $S$ 的过去有奇性 (即存在大爆炸). 此外, Theorem~\ref{thm:hawking-timelike-singularity} 蕴含这一奇性是普遍的: 满足强能量条件的膨胀 FLRW 模型的任何足够小的扰动, 都仍将有奇性. 粗略地说, 任何膨胀宇宙都必从大爆炸开始.
\item Schwarzschild 解的区域 $\{r<2m\}$ 整体双曲 [参见 Exercise~\ref{ex:causality}(\ref{ex:globally-hyperbolic-spacetime-examples})], 且满足强能量条件 (因为 $\operatorname{Ric}=0$). 度量在这一区域内可写为
\[
  g=-\mathrm{d}\tau\otimes\mathrm{d}\tau
    +\left(\frac{2m}{r}-1\right)\mathrm{d}t\otimes\mathrm{d}t
    +r^2\,\mathrm{d}\theta\otimes\mathrm{d}\theta
    +r^2\sin^2\theta\,\mathrm{d}\varphi\otimes\mathrm{d}\varphi,
\]
其中
\[
  \tau=\int_r^{2m}\left(\frac{2m}{u}-1\right)^{-1/2}\,\mathrm{d}u.
\]
因此, 可以把黑洞内部想象成一个形状随时间演化的柱面 $\mathbb{R}\times S^2$. 当 $r\to0$ 时, $S^2$ 收缩到奇点, 而 $t$ 方向伸展. 由于
\[
  \sum_{i,j=1}^3\beta_{ij}\,\mathrm{d}x^i\otimes\mathrm{d}x^j
  =\frac{\mathrm{d}r}{\mathrm{d}\tau}
   \left(-\frac{m}{r^2}\,\mathrm{d}t\otimes\mathrm{d}t+r\,\mathrm{d}\theta\otimes\mathrm{d}\theta
     +r\sin^2\theta\,\mathrm{d}\varphi\otimes\mathrm{d}\varphi\right),
\]
有
\[
  \theta=\left(\frac{2m}{r}-1\right)^{-1/2}\left(\frac{2}{r}-\frac{3m}{r^2}\right).
\]
因此, 在任意满足 $r_0<\frac{3m}{2}$ 的 Cauchy 超曲面 $S=\{r=r_0\}$ 上都有 $\theta=\theta_0<0$, 故 Theorem~\ref{thm:hawking-timelike-singularity} 保证 Schwarzschild 解在 $S$ 的未来具有奇性. 此外, Theorem~\ref{thm:hawking-timelike-singularity} 蕴含该奇性是普遍的: 满足强能量条件的 Schwarzschild 解的任意足够小的扰动, 也将有奇性. 粗略地说, 一旦坍缩进行得足够久, 就无法阻止奇点形成.
\end{enumerate}
\end{RiemannExample}

\begin{RiemannExercise}\label{ex:hawking-singularity-theorem}
\begin{enumerate}[label=(\arabic*),ref=\arabic*,leftmargin=*]
\item\label{ex:clifton-pohl-torus-incompleteness}
(\emph{Clifton–Pohl 环面}) 考虑 $\overline M=\mathbb{R}^2\setminus\{0\}$ 上的 Lorentz 度量
\[
  \overline g:=\frac{1}{u^2+v^2}(\mathrm{d}u\otimes\mathrm{d}v+\mathrm{d}v\otimes\mathrm{d}u).
\]
Lie 群 $\mathbb{Z}$ 通过
\[
  n\cdot(u,v)=(2^nu,2^nv)
\]
在 $\overline M$ 上作自由且适当的等距作用, 并在 $M=\overline M/\mathbb{Z}\cong T^2$ 上确定一个 Lorentz 度量 $g$. 证明 $(M,g)$ 不是测地完备的 (虽然 $M$ 是紧的).

(\textbf{提示}: 寻找满足 $v\equiv0$ 的零模测地线.)
\item\label{ex:anti-de-sitter-geodesic-completeness}
(\emph{二维反 de Sitter 宇宙}) 考虑带有伪 Riemann 度量
\[
  \overline g=-\mathrm{d}u\otimes\mathrm{d}u-\mathrm{d}v\otimes\mathrm{d}v
                +\mathrm{d}w\otimes\mathrm{d}w
\]
的 $\mathbb{R}^3$, 设 $(M,g)$ 是带诱导度量的子流形
\[
  H=\{(u,v,w)\in\mathbb{R}^3\mid u^2+v^2-w^2=1\}
\]
的万有覆盖. 证明:
\begin{enumerate}[label=(\alph*),leftmargin=*]
\item $(M,g)$ 的一个模型是 $M=\mathbb{R}\times(-\frac{\pi}{2},\frac{\pi}{2})$, 且
\[
  g=\frac{1}{\cos^2 x}(-\mathrm{d}t\otimes\mathrm{d}t+\mathrm{d}x\otimes\mathrm{d}x)
\]
(Fig.~\ref{fig:anti-de-sitter-exponential-map}; 因而 $(M,g)$ 并不整体双曲);
\item $(M,g)$ 是测地完备的, 但对任意 $p\in M$, $\exp_p$ 都不是满射;

(\textbf{提示}: 注意 $(\mathbb{R}^3,\overline g)$ 的每个等距映射都确定 $(M,g)$ 的一个等距映射.)
\item 存在 $p,q\in M$ 可由任意长的类时曲线连接 [参见 Exercise~\ref{ex:hawking-singularity-theorem}(\ref{ex:maximal-timelike-geodesic-between-two-events})].
\end{enumerate}
\item\label{ex:incomplete-left-invariant-lorentzian-metric}
类比第~\ref{chap:riemannian-manifolds}~章的 Exercise~\ref{ex:levi-civita-connection}(\ref{ex:hyperbolic-plane}), 可以在第~\ref{chap:differentiable-manifolds}~章 Exercise~\ref{ex:lie-groups}(\ref{ex:affine-maps-lie-group}) 的 Lie 群 $H=\mathbb{R}\times\mathbb{R}^+$ 上定义左不变 Lorentz 度量
\[
  g:=\frac{1}{x^2}(-\mathrm{d}t\otimes\mathrm{d}t+\mathrm{d}x\otimes\mathrm{d}x).
\]
证明这个度量不是测地完备的. (\textbf{注}: 这在 Riemann 几何中不可能发生 --- 参见第~\ref{chap:riemannian-manifolds}~章的 Exercise~\ref{ex:hopf-rinow-theorem}(\ref{ex:completeness-of-homogeneous-manifolds}).)
\item\label{ex:christoffel-symbols-in-synchronized-coordinates}
证明度量
\[
  g=-\mathrm{d}t\otimes\mathrm{d}t+\sum_{i,j=1}^3\gamma_{ij}\,\mathrm{d}x^i\otimes\mathrm{d}x^j
\]
的 Christoffel 符号满足
\[
  \Gamma_{00}^0=\Gamma_{00}^i=0
  \qquad\text{且}\qquad
  \Gamma_{0j}^i=\sum_{k=1}^3\gamma^{ik}\beta_{kj},
\]
其中 $(\gamma^{ij})=(\gamma_{ij})^{-1}$, $\beta_{ij}=\frac{1}{2}\frac{\partial\gamma_{ij}}{\partial t}$.
\item\label{ex:positive-pressureless-fluid-energy-contraction}
证明若 $U$ 为单位类时向量场, $V$ 为任意类时向量场, 则 $\langle U,V\rangle^2+\frac12\langle V,V\rangle$ 是正函数.
\item\label{ex:strong-energy-condition-with-cosmological-constant}
证明若时空 $(M,g)$ 的物质为静止质量密度函数为 $\rho\in C^\infty(M)$ 的无压流体, 宇宙学常数为 $\Lambda\in\mathbb{R}$ [参见 Exercise~\ref{ex:cosmology}(\ref{ex:flrw-with-cosmological-constant})], 则它满足强能量条件当且仅当 $\rho\geq0$ 且 $\rho\geq\frac{\Lambda}{4\pi}$.
\item\label{ex:locally-finite-refinements-by-simple-neighborhoods}
设 $(M,g)$ 为一个时空. 证明任意开覆盖 $\{V_\alpha\}_{\alpha\in A}$ 都有由简单邻域组成的可数局部有限加细 $\{U_n\}_{n\in\mathbb{N}}$ (即 $\bigcup_{n\in\mathbb{N}}U_n=\bigcup_{\alpha\in A}V_\alpha$, 对每个 $n\in\mathbb{N}$ 存在 $\alpha\in A$ 使 $U_n\subset V_\alpha$, 且每点 $p\in M$ 都有一个只与有限个简单邻域 $U_n$ 相交的邻域).
\item\label{ex:prove-closed-causal-futures-and-compact-diamonds}
证明 Corollary~\ref{cor:closed-causal-futures-and-compact-diamonds}.
\item\label{ex:hausdorff-closure-of-timelike-curves}
设 $(M,g)$ 为整体双曲时空, $t:M\to\mathbb{R}$ 为全局时间函数, $S=t^{-1}(0)$ 为 Cauchy 超曲面, $p\in D^+(S)$, 且 $A=D^+(S)\cap J^-(p)$. 证明在带 Hausdorff 度量的 $A$ 的所有紧子集的空间 $C(A)$ 中, 闭包 $C(S,p):=\overline{T(S,p)}$ 可以与连接 $S$ 至 $p$ 的连续因果曲线 (以 $t$ 参数化) 所成的集合等同.
\item\label{ex:maximal-timelike-geodesic-between-two-events}
设 $(M,g)$ 为整体双曲时空, $p,q\in M$ 且 $q\in I^+(p)$. 证明在所有连接 $p$ 至 $q$ 的类时曲线中, 存在一条长度最大的类时曲线, 它是类时测地线.
\item\label{ex:nonmaximal-circular-orbit-around-schwarzschild-mass}
考虑 Schwarzschild 时空中对应于半径为 $r$ 的完整一周圆形轨道的起点和终点的两个事件 $p$ 和 $q$ [参见 Exercise~\ref{ex:schwarzschild-solution}(\ref{ex:equatorial-circular-geodesics-in-schwarzschild})]. 证明相应的类时测地线不使长度达到最大.
\item\label{ex:myers-compactness-under-positive-ricci-curvature}
(\emph{Myers 定理}) 采用与 Theorem~\ref{thm:hawking-timelike-singularity} 的证明类似的思路, 证明 \textbf{Myers 定理}: 若 $(M,\langle\cdot,\cdot\rangle)$ 是完备 Riemann 流形, 且其 Ricci 曲率对某个 $\varepsilon>0$ 满足 $\operatorname{Ric}(X,X)\geq\varepsilon\langle X,X\rangle$, 则 $M$ 是紧的. 能否用这些思路证明 Riemann 几何中的一个奇性定理?
\item\label{ex:limits-of-hawking-singularity-theorem}
解释为何 Hawking 奇性定理不适用于下列各个时空:
\begin{enumerate}[label=(\alph*),leftmargin=*]
\item Minkowski 时空;
\item Einstein 宇宙 [参见 Exercise~\ref{ex:cosmology}(\ref{ex:einstein-and-de-sitter-universe-embeddings})];
\item de Sitter 宇宙 [参见 Exercise~\ref{ex:cosmology}(\ref{ex:einstein-and-de-sitter-universe-embeddings})];
\item \phantomsection\label{idx:hawking-example-anti-de-sitter}二维反 de Sitter 宇宙 [参见 Exercise~\ref{ex:hawking-singularity-theorem}(\ref{ex:anti-de-sitter-geodesic-completeness})].
\end{enumerate}
\end{enumerate}
\end{RiemannExercise}
\RGfigure[.43\textwidth]{anti-de-sitter-exponential-map}{二维反 de Sitter 宇宙中的指数映射不是满射}
